\section{RoPE 的数学表达}

\subsection{旋转矩阵的基本性质}

对位置 $m$ 的 Query 施加旋转矩阵 $R_m$，对位置 $n$ 的 Key 施加 $R_n$：

$$
\text{score}(m, n) = (R_m q)^T (R_n k) = q^T R_m^T R_n k = q^T R_{n-m} k
$$

关键性质：$R_m^T R_n = R_{n-m}$，旋转矩阵的转置等于逆旋转，两个旋转的组合只依赖位置差 $n - m$。

\begin{warningbox}{相对性是推演出来的}
我们事先只需要为每个绝对位置计算旋转矩阵，相对位置编码是数学推导的自然结果，不需要显式构造。
\end{warningbox}

\subsection{频率参数 $\theta$}

RoPE 在 embedding 向量的不同 2D 子空间（subspace）上施加不同频率的旋转：

$$
\theta_i = \text{base}^{-2i/d}
$$

\begin{itemize}
  \item $\text{base}$：一般取 10000（标准 RoPE），Qwen3 取 \textbf{1000000}（追求更大 context window）
  \item $i$：子空间索引，从 0 到 $d/2 - 1$
  \item $d$：head dimension
\end{itemize}

$\theta_i$ 随 $i$ 增大呈\textbf{指数级衰减}，取 $\log$ 后呈线性下降。

\subsection{本章小结}

RoPE 通过绝对位置的旋转矩阵实现相对位置编码，$\theta_i$ 控制了各子空间的旋转频率。base 越大，低频子空间的旋转越慢，有利于长距离建模。

